\documentclass[11pt]{article}
\usepackage[a4paper,margin=1in]{geometry}
\usepackage{amsmath,amssymb,amsthm}
\usepackage{hyperref}

\theoremstyle{plain}
\newtheorem{theorem}{Theorem}[section]
\newtheorem{lemma}[theorem]{Lemma}
\newtheorem{proposition}[theorem]{Proposition}
\theoremstyle{definition}
\newtheorem{definition}[theorem]{Definition}
\theoremstyle{remark}
\newtheorem{remark}[theorem]{Remark}

% ------------------------------------------------------------------
% Notation, following Shoenfield (Mathematical Logic, Ch.6) as used in
% the companion paper arGII, with code(.) for the Goedel encoding.
% ------------------------------------------------------------------
\newcommand{\code}[1]{\mathit{code}(#1)}
\newcommand{\thmT}{\mathbf{th}}
\newcommand{\num}{\mathbf{num}}
\newcommand{\evalU}{\mathbf{U}}
\newcommand{\parse}{\mathbf{parse}}
\newcommand{\stepU}{\mathbf{step}}
\newcommand{\defp}{\mathbf{def}}
\newcommand{\CK}{\mathbf{CK}}
\newcommand{\Kc}{K}
\newcommand{\Fun}{\mathit{Fun}}
\newcommand{\Term}{\mathit{Term}}
\newcommand{\Con}{\mathrm{Con}}
\newcommand{\Thm}[1]{\mathit{Thm}(#1)}
\newcommand{\fuelG}{\mathbf{fuel}}
\newcommand{\fuelMu}{\mathbf{fuel}_\mu}
\newcommand{\Lstar}{L^{\!\star}}
\newcommand{\rta}{\rightarrow}
\newcommand{\imp}{\supset}

\title{A surprise-examination (sorites) proof of G\"odel's second\\
  incompleteness theorem in Basic Recursive Arithmetic}
\author{Thierry Coquand}
\date{\today}

\begin{document}
\maketitle

\begin{abstract}
We describe a second, independent machine-checked proof of G\"odel's
second incompleteness theorem for Church's Basic Recursive Arithmetic
(BRA), following Chaitin's information-theoretic route and the
surprise-examination argument of Kritchman and Raz. The formal system,
its numerals, the encoding $\code{\cdot}$, the verifier $\thmT$ and the
numeral functor $\num$ are exactly those of the companion paper
\cite{arGII}, which gives a first proof of the same theorem along
Guard's lines; we do not repeat their construction here. The two
ingredients specific to this proof are (i)~an \emph{internal} form of
Chaitin's first incompleteness theorem --- a single BRA implication
built around an object-level interpreter --- and (ii)~a \emph{sorites},
or heap-of-sand, descent that turns Chaitin's theorem into the second
incompleteness theorem. The whole development is checked in Agda under
\texttt{-{}-safe -{}-without-K -{}-exact-split} with no postulates and no
holes, from the single hypothesis that $T$ does not prove $0=1$.
\end{abstract}

\section{The statement}

We work entirely inside the system $T$ of Basic Recursive Arithmetic in
the formulation of Guard, with Shoenfield's notation; the axioms, the
numerals $\mathbf{k}_n = \mathbf{s}^n(\mathbf{O})$, the encoding
$\code{\cdot}$, the verifier $\thmT \in \Fun_1$ (with
$\Thm{\thmT(\code{d}) = \code{A}}$ whenever $d$ codes a derivation of
$A$), the numeral functor $\num \in \Fun_1$ (with
$\Thm{\num(\mathbf{k}_n) = \code{\mathbf{k}_n}}$) and the
derivability-internalising Theorems~12 and~13 are all as set up in the
companion paper~\cite{arGII}, to which we refer for every detail of the
system. We write $\Thm{A}$ for ``$A$ is a
theorem of $T$''.

As in~\cite{arGII}, the consistency of $T$ is the \emph{open} formula
\[
  \Con_T \;\equiv\; \neg\,(\thmT(\mathbf{x}_0) = \code{\mathbf{O}=\mathbf{s}\,\mathbf{O}}),
\]
with the single free variable $\mathbf{x}_0$. Since BRA has no object
quantifiers, a theorem is an open formula read as implicitly universally
quantified over its free variables, so $\Con_T$ says: \emph{for every
term $\mathbf{x}_0$, $\mathbf{x}_0$ does not code a derivation of}
$\mathbf{O}=\mathbf{s}\,\mathbf{O}$. The theorem we prove is the same
meta-implication that the Guard-route proof establishes,
\[
  \boxed{\;\Thm{\Con_T} \;\Longrightarrow\; \Thm{\mathbf{O}=\mathbf{s}\,\mathbf{O}},\;}
\]
i.e.\ if $T$ proves its own consistency then $T$ is inconsistent. The
arrow $\Longrightarrow$ is a meta-level implication between derivability
judgements, realised in Agda as a total function transforming any proof
of $\Con_T$ into a proof of $\mathbf{O}=\mathbf{s}\,\mathbf{O}$.
The consistency hypothesis is consumed exactly once per step of the
descent of \S\ref{sec:sorites}, and is the \emph{only} assumption: every
other ingredient is constructed.

\section{Chaitin's first incompleteness theorem, internalised}
\label{sec:chaitin}

Chaitin's proof of the first incompleteness theorem~\cite{Chaitin71,Chaitin74}
replaces the liar sentence by Berry's paradox and Kolmogorov complexity.
We need this argument not as a metatheorem but as a single BRA
implication, with the underlying universal machine realised as an honest
object-level program. This section describes that realisation in
ordinary computational terms; the detailed Agda construction is part of
the development cited in \S\ref{sec:formal}.

\subsection{Programs as ternary numbers}

A \emph{program} is simply a natural number. Read a number $p$ in base
$3$; its digits, drawn from the three-symbol alphabet $\{1,2,3\}$, form a
string, and that string is fed to a fixed parser $\parse$ that turns it
into a syntax tree --- an expression of the object language. The
\emph{length} $|p|$ of the program is the number of base-$3$ digits, so
\[
  |p| \le n \quad\Longleftrightarrow\quad p < 3^{\,n+1}.
\]
Thus the enumeration of programs is the \emph{identity}: the $k$-th
program is the number $k$, and the finite set of programs of length at
most $\Lstar$ is exactly the initial segment
\[
  \{\,p : p < N\,\}, \qquad N := 3^{\,\Lstar+1}.
\]
This is the single simplification that makes the whole argument
feasible: there is no separate enumeration to verify --- ``the set of
short programs'' \emph{is} an interval of numbers, and quantification
over it is bounded quantification over $\{p<N\}$.

\subsection{The universal interpreter as a CK machine}

The interpreter $\evalU$ that runs a parsed program is realised as a small
abstract machine of the Landin CK kind. A \emph{configuration} is a pair
\[
  \langle\, c,\; \kappa \,\rangle,
\]
where $c$ is the expression currently under evaluation (the
\emph{control}) and $\kappa$ is a \emph{continuation stack} recording the
contexts still to be returned to. A one-step transition function $\stepU$
fires a single reduction rule: it either decomposes $c$, pushing the
deferred work onto $\kappa$, or, when $c$ has reached a value, pops the
top frame of $\kappa$ and plugs the value into it. Running the program
for $l$ steps is the $l$-fold iteration $\stepU^{\,l}$ started from the
initial configuration; the machine halts when the control is a value and
the stack is empty, and the output is read off that final configuration.
Around this evaluator we wrap one more layer, a bounded
$\mu$-search loop, so that the outer behaviour of $\evalU$ matches the
informal recipe ``enumerate $\thmT(0), \thmT(1), \dots$ until a fitting
proof is found'' used by the diagonal program below.

For a program $p$ we write
\[
  \defp_p(l, u) \;\equiv\; \text{``running $p$ for $l$ steps halts with output $u+1$''}
\]
(the $+1$ separates ``halted with value $u$'' from ``not yet halted'',
coded by $0$); $\defp_p \in \Fun_2$ for each $p$. Folding the bounded
disjunction of the $\defp_p$ over the finite set $\{p<N\}$ gives a single
object function $\CK \in \Fun_2$ with
\[
  \CK(u, x) = \mathbf{O}
  \;\Longleftrightarrow\;
  \text{some program $p<N$, run for $x$ steps, outputs $u+1$,}
\]
i.e.\ $\CK(u,x)=\mathbf{O}$ expresses $\Kc(u)\le\Lstar$ at run-length
$x$, and the Kolmogorov bound
\[
  \Kc(u) > \Lstar \quad\text{is expressed by}\quad \CK(u,x)\neq\mathbf{O}
\]
with $x$ a free run-length variable.

\subsection{Term-valued fuel: a Gandy--Howard majorant}
\label{sec:fuel}

The delicate point is not the diagonal argument --- which is a faithful
rendering of Chaitin's --- but the fact that the interpreter must be run
\emph{inside} a $\thmT$-fact. Since $\thmT$, $\evalU$ and $\stepU$ all
take object terms, the number of steps appearing in such a fact must
itself be an object \emph{term}, not a meta-level natural; and
Theorem~13 of~\cite{arGII} internalises only statements about BRA terms.
A completeness statement of the naive form ``for every numeral $\bar n$
large enough, $\stepU^{\,\bar n}$ on the program halts'' is therefore
useless: it is quantified over meta-naturals and cannot be fed into an
internal derivation. What is needed instead is a completeness statement
with a closed \emph{combinator} in the fuel slot,
\[
  (x:\Term)\to\;\mathit{Thm}\,
     \bigl(\,\stepU^{\,\fuelG(f)(x)}(\text{init}(f,x)) = \text{halt}(f\cdot x)\,\bigr),
\]
where $\fuelG(f)\in\Fun_1$ is built once and for all from $f$.

Proving such a statement by induction on the combinator tree of $f$
stalls at the primitive-recursion node, where one would need numerical
control over the running time of the recursive calls --- information not
available at the derivation level. The resolution adapts the classical
\emph{majorisation} method of Howard and Gandy~\cite{Howard73,Gandy80}:
one defines, by primitive recursion mirroring the structure of $f$, a
closed combinator $\fuelG(f)$ that provably \emph{majorises} the actual
number of steps any run of $f$ consumes at every term input. The
induction then carries the stable shape ``the run completes at
\emph{any} fuel $\ge \fuelG(f)(x)$'', which descends through the
recursion node because $\fuelG$ is itself recursive in step with $f$. The
$\mu$-search loop needs the same idea one level higher, with a majorant
$\fuelMu$ stacked on top of $\fuelG$ by a fuel-composition lemma. The
closed term that finally occupies the fuel slot is
$\fuelMu(\cdots) + \fuelG(\cdots)$. Without this majorisation there would
simply be no term to put there.

\subsection{The diagonal program and the internal implication}

Chaitin's short program $g_{\Lstar}$ enumerates $\thmT(0), \thmT(1),
\dots$ until it meets the code of a statement of the form
$\Kc(?)>\Lstar$, and then outputs the witnessed number $?$. Its size is a
constant plus the number of digits needed to write $\Lstar$, that is
$c + \log\Lstar$, which is below $\Lstar$ once $\Lstar$ is large. The
program $g_{\Lstar}$ is obtained as a genuine one-variable object
combinator $G\in\Fun_1$ by bracket abstraction --- an honest BRA program,
not a meta-level term transformer --- so that $G\cdot w$ is literally the
constructed code, as a \emph{proved} $T$-equation. Let $\mathbf{out}(w)$
be the object function that reads the witness off $w$. The internalised
Chaitin theorem is the single BRA implication
\[
  \mathit{Thm}\Bigl(\;\bigl(\thmT(w) = \code{\,\Kc(\mathbf{out}(w))>\Lstar\,}\bigr)
        \;\imp\;
        \bigl(\thmT(G\cdot w) = \code{\mathbf{O}=\mathbf{s}\,\mathbf{O}}\bigr)\;\Bigr),
  \tag{C}
\]
valid for \emph{every} term $w$, with no consistency premise. In words:
$T$ proves that \emph{if} $\thmT(w)$ is the incompressibility statement
about $w$'s own read-back, \emph{then} the code $G\cdot w$ is a $T$-proof
of $0=1$. The deduction lives entirely inside the object logic as an
implication; the run of the interpreter on $g_{\Lstar}$, with the
Term-fuel of \S\ref{sec:fuel}, is what closes the positive leg of~(C).
This is the only place where the evaluator machinery is needed; the
descent of the next section uses~(C) purely as a black box.

\section{The sorites descent}
\label{sec:sorites}

Chaitin's theorem says that, for a large enough $\Lstar$, no statement
$\Kc(x)>\Lstar$ is provable in a consistent $T$, even though for each
$x$ exactly one of $\Kc(x)\le\Lstar$, $\Kc(x)>\Lstar$ holds and the
former, when true, is provable (by exhibiting and running the short
program). Kritchman and Raz~\cite{KritchmanRaz} observed that this
tension can be wound into the second incompleteness theorem by an
argument modelled on the surprise-examination paradox.

\subsection{The counting argument}

Fix $\Lstar$ as in \S\ref{sec:chaitin} and put $N = 3^{\Lstar+1}$. Let
\[
  m \;=\; \#\{\,u \in [0,N] : \Kc(u) > \Lstar\,\}.
\]
Because there are only $N$ programs of length $\le\Lstar$, at most $N$ of
the $N+1$ values in $[0,N]$ can be compressible, so at least one is
incompressible: $m \ge 1$, and this is provable in $T$. Kritchman and Raz
then prove, \emph{by induction on $i$ and inside $T$}, that
$\Thm{m \ge i}$ for $i = 1, 2, \dots$. The induction step
$\Thm{m\ge i}\Rightarrow\Thm{m\ge i{+}1}$ is exactly where Chaitin's
theorem and one use of consistency enter (it is their displayed
equations~(1) for Chaitin and~(2) for the $\Sigma_1$-completeness of
$\Kc(\cdot)\le\Lstar$). But $m \le N+1$ outright, so the induction cannot
run forever: once $i$ passes $N+1$, the provable statement $m \ge i$
contradicts $m \le N+1$, and $T$ proves $0=1$.

\subsection{Why this is a sorites, not a vicious circle}

The original surprise-examination paradox (``there will be an exam next
week but you will not know the day until it is held'') is usually
analysed as a self-referential knowledge puzzle. The mathematical core
that survives formalisation, however, has the shape of the
\emph{sorites}, or heap-of-sand, paradox. Think of the interval $[0,N]$
--- equivalently the population of incompressible numbers counted by $m$
--- as a heap of sand. Each step of the induction removes one grain: it
is, taken on its own, an entirely sound BRA derivation, consisting of one
internal application of the Chaitin implication~(C) reflected through a
single use of consistency. No individual grain-removal is the culprit.
Yet after $N+1$ removals the heap is gone and $T$ has proved
$\mathbf{O}=\mathbf{s}\,\mathbf{O}$. The contradiction lives only in the
iteration, exactly as in the heap paradox, and it is the second
incompleteness theorem --- the impossibility of $T$ proving its own
consistency --- that explains why the seemingly innocuous step cannot in
fact be taken all the way (cf.\ the resolution discussed
in~\cite{KritchmanRaz}: the missing premise is the consistency of $T$
itself).

\subsection{The formal realisation: six steps}
\label{sec:sixsteps}

In the formalisation the count $m$ is carried by an open formula
$S(r)$, indexed by a stage $r \in [0,N]$ and built as a negated big
conjunction over ``days''. For programs $p_r,\dots,p_N$ (each of length
$\le\Lstar$, ranging over $\{p<N\}$) write
\[
  \Kc(\mathbf{x}_0; p_r,\dots,p_N) \;\equiv\;
    \bigwedge_{d=r}^{N} \defp_{p_d}(\mathbf{x}_0, d),
\]
``each program $p_d$ describes the day $d$, at the common run-length
$\mathbf{x}_0$''. The stage formula is
\[
  S(r) \;\equiv\;
    \Thm{\neg\,\Kc(\mathbf{x}_0; p_r,\dots,p_N)}
    \qquad\text{(for every choice of $p_r$).}
\]
The descent climbs $S(0) \to S(1) \to \dots \to S(N)$, one grain at a
time, and the inductive step $S(r)\to S(r+1)$ is the six-step block that
realises Kritchman--Raz's induction step; it follows the high-level note
\texttt{clos} of the development and reuses the implication~(C) as a
black box.

\begin{description}
\item[Step 1 (run-monotonicity and the finite set).]
From $S(r)$ one derives the implication
``$\Kc(\mathbf{x}_0; p_{r+1},\dots,p_N) \imp \neg\defp_{p_r}(\mathbf{x}_1,r)$''
at an \emph{independent} run-length $\mathbf{x}_1$, the two fuels
$\mathbf{x}_0,\mathbf{x}_1$ being aligned to a common bound by
monotonicity of the interpreter in its step count (running longer never
unsays a halt). Folding the bounded disjunction over the finite set
$\{p<N\}$ rewrites the consequent as $\Kc(r) > \Lstar$, giving
\[
  \mathit{Thm}\bigl(\,\Kc(\mathbf{x}_0; p_{r+1},\dots,p_N) \imp (\Kc(r)>\Lstar)\,\bigr).
\]

\item[Step 2 (encode and substitute the numeral).]
By Theorem~13 of~\cite{arGII} the Step-1 implication is internalised: one
obtains a term $f$ with
$\Thm{\thmT(f\cdot\mathbf{x}_0)
   = \code{\,\Kc(\num\,\mathbf{x}_0; p_{r+1},\dots,p_N) \imp \Kc(r)>\Lstar\,}}$,
the subject slot $\mathbf{x}_0$ being replaced inside the code by its
numeral $\num\,\mathbf{x}_0$ (the consequent has no free $\mathbf{x}_0$
and is left inert).

\item[Step 3 (internal modus ponens).]
An encoded modus ponens (the Carneiro-lifted $\mathbf{imp\_encoded\_mp}$
of~\cite{arGII}) chains Step~2 with a $\thmT$-fact for the antecedent.

\item[Step 4 (Theorem~13 for the conjunction).]
Writing the big conjunction as $K_r(\mathbf{x}_0)=\mathbf{O}$, the
$\Sigma_1$-completeness direction
$\bigl(K_r(\mathbf{x}_0)=\mathbf{O}\bigr) \imp
  \thmT(D\,\mathbf{x}_0) = \code{\,\Kc(\num\,\mathbf{x}_0; p_{r+1},\dots,p_N)\,}$
supplies that antecedent fact (Kritchman--Raz's equation~(2)).

\item[Step 5 (Chaitin).]
The internal Chaitin implication~(C) is applied:
$\mathit{Thm}\bigl(\,(\thmT(w)=\code{\Kc(r)>\Lstar})
   \imp (\thmT(G\cdot w)=\code{\mathbf{O}=\mathbf{s}\,\mathbf{O}})\,\bigr)$.
Composing Steps~3--5 yields
$\mathit{Thm}\bigl(\,K_r(\mathbf{x}_0)=\mathbf{O}
   \imp (\thmT(h\cdot\mathbf{x}_0)=\code{\mathbf{O}=\mathbf{s}\,\mathbf{O}})\,\bigr)$.

\item[Step 6 (consistency).]
Finally $\Con_T$ is instantiated at $h\cdot\mathbf{x}_0$ and, by
contraposition, refutes the consequent --- so the antecedent fails:
$\Thm{\neg\,\Kc(\mathbf{x}_0; p_{r+1},\dots,p_N)}$, which is $S(r+1)$.
This is the single use of consistency per grain (Kritchman--Raz's
equation~(1)).
\end{description}

\paragraph{Base case.}
$S(0)$ is hypothesis-free and combinatorial. There are $N+1$ days but
only $N$ program slots, so any assignment of a program to each day forces
a collision $p_i = p_j$ with $i\neq j$; the two day-conjuncts then assert
that one program, at the same run-length, outputs both $\mathbf{s}^i$ and
$\mathbf{s}^j$, which is impossible. Hence $\neg\Kc(\mathbf{x}_0;p_0,\dots,p_N)$.

\paragraph{Final clash.}
At the top stage $r=N$ the remaining conjunction is empty, hence
provably true, so the Step-1 implication fires \emph{unconditionally}:
$T$ proves $\Kc(N)>\Lstar$, the Chaitin diagonal makes its provability
into a proof of $0=1$, and $\Con_T$ refutes it. The heap is empty and
\[
  \Thm{\mathbf{O}=\mathbf{s}\,\mathbf{O}}.
\]
By external induction the chain $S(0)\to\dots\to S(N)$ closes the proof.
This induction is not informal: it is carried out in the metatheory by
Agda, as a recursion on the meta-natural $r$ that turns the six-step
block into a function
$\Thm{S(r)} \to \Thm{S(r+1)}$ and iterates it from
the base derivation of $S(0)$ up to $S(N)$.

\section{The formalisation}
\label{sec:formal}

The proof is checked in Agda, in the same \texttt{T4} development as the
Guard-route proof of~\cite{arGII}, under the global options
\texttt{-{}-safe -{}-without-K -{}-exact-split}, with no postulates and no
holes. The headline is a single function from a derivation of $\Con_T$ to
a derivation of $\mathbf{O}=\mathbf{s}\,\mathbf{O}$; $\Con_T$ is the only
hypothesis, and it is the same open consistency formula as
in~\cite{arGII}.

The one genuinely non-classical piece is the Gandy--Howard majorant of
\S\ref{sec:fuel}. The diagonalisation is Chaitin's, but the closed object
term occupying the fuel slot of the internal interpreter run has no
counterpart in the informal argument, and its construction (by recursion
mirroring the program, with a substituted-motive induction through the
recursion node) is what lets evaluator completeness be \emph{stated and
proved} inside BRA. It is worth noting, finally, that Guard's
formulation of Basic Recursive Arithmetic --- syntax-first, with the
verifier $\thmT$ and the numeral functor $\num$ as honest object
functions --- proved well suited to expressing all of these results,
including the internal Chaitin implication and the surprise descent.

The development, including both proofs of the second incompleteness
theorem, is available at
\url{https://github.com/coquand/agda-godel-tree}.

\begin{thebibliography}{10}

\bibitem{arGII}
T.~Coquand.
\newblock Autoformalisation of G\"odel's second incompleteness theorem in
  Basic Recursive Arithmetic.
\newblock arXiv:2606.01898, 2026.
\newblock \url{https://arxiv.org/abs/2606.01898}.

\bibitem{Chaitin71}
G.~J. Chaitin.
\newblock Computational complexity and G\"odel's incompleteness theorem.
\newblock {\em ACM SIGACT News}, 9:11--12, 1971.

\bibitem{Chaitin74}
G.~J. Chaitin.
\newblock Information-theoretic limitations of formal systems.
\newblock {\em Journal of the ACM}, 21(3):403--424, 1974.

\bibitem{Gandy80}
R.~O. Gandy.
\newblock Proofs of strong normalization.
\newblock In J.~P. Seldin and J.~R. Hindley, editors, {\em To H.~B. Curry:
  Essays on Combinatory Logic, Lambda Calculus and Formalism}, pages
  457--477. Academic Press, 1980.

\bibitem{Guard63}
J.~R. Guard.
\newblock Lecture notes on recursive arithmetic.
\newblock Unpublished lecture notes, 1963.

\bibitem{Howard73}
W.~A. Howard.
\newblock Hereditarily majorizable functionals of finite type.
\newblock In A.~S. Troelstra, editor, {\em Metamathematical Investigations of
  Intuitionistic Arithmetic and Analysis}, volume 344 of {\em Lecture Notes in
  Mathematics}, pages 454--461. Springer, 1973.

\bibitem{Kikuchi97}
M.~Kikuchi.
\newblock Kolmogorov complexity and the second incompleteness theorem.
\newblock {\em Archive for Mathematical Logic}, 36:437--443, 1997.

\bibitem{KritchmanRaz}
S.~Kritchman and R.~Raz.
\newblock The surprise examination paradox and the second incompleteness
  theorem.
\newblock {\em Notices of the American Mathematical Society}, 57(11):1454--1458,
  2010.

\bibitem{Shoenfield}
J.~R. Shoenfield.
\newblock {\em Mathematical Logic}.
\newblock Addison--Wesley, Reading, MA, 1967.

\end{thebibliography}

\end{document}
